\documentclass[11pt,a4paper]{article}

\usepackage[utf8]{inputenc}
\usepackage[T1]{fontenc}
\usepackage[margin=1in]{geometry}
\usepackage{amsmath,amssymb,amsthm,amsfonts}
\usepackage{mathtools}
\setcounter{MaxMatrixCols}{20}
\usepackage{bm}
\usepackage{graphicx}
\usepackage{xcolor}
\usepackage{hyperref}
\usepackage{enumitem}
\usepackage{booktabs}
\usepackage{tikz}
\usepackage{listings}
\usepackage{caption}
\usepackage{titling}
\usepackage{fancyhdr}
\usepackage{parskip}
\usepackage{tcolorbox}
\tcbuselibrary{breakable}
\usepackage{longtable}

\usetikzlibrary{arrows.meta,positioning,shapes.geometric,calc}

\definecolor{codebg}{RGB}{245,245,245}
\definecolor{codekw}{RGB}{0,0,180}
\definecolor{codecmt}{RGB}{0,128,0}
\definecolor{codestr}{RGB}{180,0,0}

\lstdefinestyle{mystyle}{
    backgroundcolor=\color{codebg},
    commentstyle=\color{codecmt}\itshape,
    keywordstyle=\color{codekw}\bfseries,
    stringstyle=\color{codestr},
    basicstyle=\ttfamily\footnotesize,
    breakatwhitespace=false,
    breaklines=true,
    keepspaces=true,
    numbers=left,
    numbersep=5pt,
    showspaces=false,
    showstringspaces=false,
    showtabs=false,
    tabsize=2,
    frame=single,
    framerule=0.4pt,
    rulecolor=\color{black!30}
}
\lstset{style=mystyle}

\hypersetup{
    colorlinks=true,
    linkcolor=blue!60!black,
    citecolor=green!60!black,
    urlcolor=blue!60!black,
    pdftitle={EKF Parameters: Theory to Practice},
    pdfauthor={Multi-Sensor Fusion Project}
}

\newtcolorbox{paramcard}[1]{colback=blue!3!white,colframe=blue!60!black,fonttitle=\bfseries\ttfamily,title=#1,breakable}
\newtcolorbox{warnbox}[1]{colback=red!5!white,colframe=red!75!black,fonttitle=\bfseries,title=#1,breakable}
\newtcolorbox{insightbox}[1]{colback=green!5!white,colframe=green!50!black,fonttitle=\bfseries,title=#1,breakable}
\newtcolorbox{intuitionbox}[1]{colback=yellow!5!white,colframe=orange!75!black,fonttitle=\bfseries,title=#1,breakable}
\newtcolorbox{mathbox}[1]{colback=violet!3!white,colframe=violet!60!black,fonttitle=\bfseries,title=#1,breakable}

\pagestyle{fancy}
\fancyhf{}
\fancyhead[L]{\small EKF Parameters --- Theory to Practice}
\fancyhead[R]{\small Husarion Lynx Project}
\fancyfoot[C]{\thepage}

\title{\Huge\bfseries EKF Parameters\\[0.4em]\Large Theory $\leftrightarrow$ Practice\\[0.3em]\large Each Knob Explained:\\Intuition, Math, Effect\\[0.5em]\normalsize\itshape Husarion Lynx --- \texttt{robot\_localization} dual-EKF + NavSat}
\author{Multi-Sensor Fusion Project}
\date{\today}

\begin{document}

\maketitle

\begin{abstract}
\noindent Companion guide to the EKF Deep Dive document. Where the deep dive built theory bottom-up, this document goes parameter-by-parameter through the actual \texttt{robot\_localization} configuration used in the Husarion Lynx project (\texttt{dual\_ekf\_navsat.yaml}), explaining for every knob: (i) what the parameter does mechanically inside the EKF; (ii) the mathematical equation it appears in; (iii) the intuition --- what increasing or decreasing it physically means; (iv) the effect on filter behavior; (v) typical failure modes when misconfigured; (vi) how to tune it. Notation throughout follows Thrun-Burgard-Fox: state mean $\mu_t$, covariance $\Sigma_t$, motion noise covariance $R_t$, measurement noise covariance $Q_t$. Note: the \texttt{robot\_localization} YAML uses the engineering convention where the YAML key \texttt{process\_noise\_covariance} stores what we call $R_t$ in Thrun notation.
\end{abstract}

\tableofcontents
\newpage

% ============================================================
\section{The Three EKFs in this Project}
% ============================================================

The Husarion Lynx setup runs three EKF instances plus one transform node. Understanding which one each parameter belongs to is the first step.

\begin{center}
\footnotesize
\begin{tabular}{p{3.2cm} p{2.0cm} p{4.5cm} p{3.5cm}}
\toprule
\textbf{Node} & \textbf{Frame} & \textbf{Inputs} & \textbf{Outputs} \\
\midrule
\texttt{ekf\_filter\_node\_odom} (EKF \#1, Local) & \texttt{world\_frame: odom} & wheel + IMU + VO + LiDAR + control & \texttt{/odometry/filtered}, TF \texttt{odom}$\to$\texttt{base\_link} \\
\texttt{ekf\_filter\_node\_map} (EKF \#2, Global) & \texttt{world\_frame: map} & wheel + IMU + VO + LiDAR (\textbf{NO GPS}) & \texttt{/odometry/filtered\_map}, TF \texttt{map}$\to$\texttt{odom} \\
\texttt{navsat\_transform} & --- & GPS (\texttt{NavSatFix}) + IMU + \texttt{/odometry/filtered} & \texttt{/odometry/gps} (analysis only) \\
\bottomrule
\end{tabular}
\end{center}

\begin{warnbox}{Critical Architectural Decision}
This project deliberately departs from the canonical \texttt{robot\_localization} dual-EKF + NavSat tutorial. \textbf{GPS is NOT fused into EKF \#2.} Both EKFs receive identical sensor inputs (wheel + IMU + VO + LiDAR), differing only in \texttt{world\_frame}.

\textbf{Reason}: GPS noise injected through cross-covariance terms $\Sigma_{x,\psi}$ and $\Sigma_{y,\psi}$ caused sustained map frame yaw drift of $\sim 3.56°/$s in earlier configurations. Without a strong absolute yaw observation (no magnetometer in this Gazebo IMU, no dual-antenna RTK), nothing in the filter can fight the leak. The covariance cross-term problem is the single most important reason to read this document carefully before re-enabling GPS in EKF \#2.

GPS still flows through the pipeline (\texttt{navsat\_transform\_node} runs, \texttt{/odometry/gps} is published) but neither EKF subscribes to that topic. GPS is available for offline analysis notebooks comparing filter output to GPS truth, but does \emph{not} influence the live state estimate.

See Section~\ref{sec:gps-exclusion} for the full mathematical derivation, numerical example, trade-offs, and alternatives.
\end{warnbox}

\textbf{Why this dual-EKF then?} With both EKFs receiving identical inputs, the resulting estimates converge to nearly the same value, so $T_{\text{map}}^{\text{odom}} \approx I$. EKF \#1 publishes \texttt{odom}$\to$\texttt{base\_link} for the controller (smooth, continuous). EKF \#2 publishes \texttt{map}$\to$\texttt{odom} (near-identity, also smooth in this setup). The map frame is \emph{anchored at the robot's spawn position} via sensor consensus, not GPS.

\textbf{Sensor noise measured} for this hardware (from \texttt{plots/summary.md}):
\begin{itemize}
\item GPS: $\sigma_{\text{lat}}^2 = 0.263$ m$^2$, $\sigma_{\text{lon}}^2 = 0.245$ m$^2$, $\sigma_{\text{alt}}^2 = 1.230$ m$^2$
\item IMU gyro: $\sigma_x^2 = 2.26 \times 10^{-4}$, $\sigma_y^2 = 1.46 \times 10^{-3}$, $\sigma_z^2 = 4.16 \times 10^{-2}$ (rad/s)$^2$
\item IMU accel: $\sigma_x^2 = 0.97$, $\sigma_y^2 = 0.57$, $\sigma_z^2 = 0.69$ (m/s$^2$)$^2$
\item Wheel odom: pose\_cov $= 0.001$, twist\_cov $= 0.001$ (from \texttt{diff\_drive\_controller})
\item Map frame stability (current config, GPS excluded): $\pm 0.1$ m, $\pm 1°$ over 2 minutes
\end{itemize}

\textbf{Notation reminder.} In the YAML keys, \texttt{process\_noise\_covariance} corresponds to Thrun's $R_t$ (motion uncertainty). The covariance fields inside each \texttt{nav\_msgs/Odometry} or \texttt{sensor\_msgs/Imu} message correspond to Thrun's $Q_t$ (measurement uncertainty).

% ============================================================
\section{State Vector Refresher}
% ============================================================

Every \texttt{ekf\_node} in \texttt{robot\_localization} maintains a 15-dimensional state:

\[
\mu_t = [\,\underbrace{x, y, z}_{\text{position}},\, \underbrace{\phi, \theta, \psi}_{\text{roll, pitch, yaw}},\, \underbrace{\dot{x}, \dot{y}, \dot{z}}_{\text{linear vel}},\, \underbrace{\dot{\phi}, \dot{\theta}, \dot{\psi}}_{\text{angular vel}},\, \underbrace{\ddot{x}, \ddot{y}, \ddot{z}}_{\text{linear accel}}\,]^\top
\]

Every parameter that takes a 15-element list (e.g., \texttt{process\_noise\_covariance} diagonal, \texttt{initial\_estimate\_covariance} diagonal) maps element-wise to this ordering. Memorise the order: position, RPY, linear velocities, angular velocities, linear accelerations.

\textbf{Index reference:}

\begin{center}
\footnotesize
\begin{tabular}{cc|cc|cc}
\toprule
Idx & Component & Idx & Component & Idx & Component \\
\midrule
0 & $x$ & 5 & $\psi$ (yaw) & 10 & $\dot{\theta}$ \\
1 & $y$ & 6 & $\dot{x}$ (vx) & 11 & $\dot{\psi}$ (vyaw) \\
2 & $z$ & 7 & $\dot{y}$ (vy) & 12 & $\ddot{x}$ (ax) \\
3 & $\phi$ (roll) & 8 & $\dot{z}$ (vz) & 13 & $\ddot{y}$ (ay) \\
4 & $\theta$ (pitch) & 9 & $\dot{\phi}$ (vroll) & 14 & $\ddot{z}$ (az) \\
\bottomrule
\end{tabular}
\end{center}

% ============================================================
\section{Frame Parameters}
% ============================================================

These define the coordinate frames in the TF tree.

\begin{paramcard}{map\_frame}
\textbf{Type:} string \quad \textbf{Project value:} \texttt{map}

\textbf{What:} Name of the global, world-fixed reference frame.

\textbf{Math role:} The frame in which $\mu_t$ is expressed when \texttt{world\_frame: map}. All transforms involving GPS, ENU coordinates terminate here.

\textbf{Intuition:} ``The unmoving big-picture frame.'' If you stood outside the robot with a god-view, this is your frame. Origin is set by GPS datum or pre-defined.

\textbf{Effect:} Only meaningful for the \emph{global} EKF (EKF \#2). Local EKF still has it set but ignores it (since \texttt{world\_frame: odom}). Must match across all nodes that share a TF tree.

\textbf{Failure mode:} Misnaming silently disconnects the TF tree --- planner finds no path to base\_link.
\end{paramcard}

\begin{paramcard}{odom\_frame}
\textbf{Type:} string \quad \textbf{Project value:} \texttt{odom}

\textbf{What:} Name of the locally drifting world-fixed reference frame.

\textbf{Math role:} The frame for $\mu_t$ when \texttt{world\_frame: odom}. The local EKF (EKF \#1, EKF \#3) operates here.

\textbf{Intuition:} ``Where the robot thinks it is, ignoring GPS.'' Drifts with time but is locally smooth and continuous --- exactly what the controller needs.

\textbf{Effect:} The local EKF publishes \texttt{odom}$\to$\texttt{base\_link}. The global EKF publishes \texttt{map}$\to$\texttt{odom}. The composition gives global pose.

\textbf{Failure mode:} If two nodes claim to publish \texttt{odom}$\to$\texttt{base\_link}, TF tree breaks (tf complains about multiple parents).
\end{paramcard}

\begin{paramcard}{base\_link\_frame}
\textbf{Type:} string \quad \textbf{Project value:} \texttt{base\_link}

\textbf{What:} Frame rigidly attached to robot body.

\textbf{Math role:} All sensor measurements are transformed into this frame using static transforms from URDF before fusion. The motion model $g(u_t, \mu_{t-1})$ predicts the pose of \texttt{base\_link} relative to the world frame.

\textbf{Intuition:} ``The robot itself.'' Convention: $+x$ forward, $+y$ left, $+z$ up (right-hand rule).

\textbf{Effect:} All input sensor messages must have their \texttt{frame\_id} resolvable to \texttt{base\_link} via static TF. Wrong sensor frame $\to$ measurement applied at wrong rotation/translation $\to$ filter biases.
\end{paramcard}

\begin{paramcard}{world\_frame}
\textbf{Type:} string \quad \textbf{Project value:} \texttt{odom} (EKF \#1, \#3) / \texttt{map} (EKF \#2)

\textbf{What:} The frame in which the EKF maintains its state estimate.

\textbf{Math role:} Determines which TF transform the node publishes:
\begin{itemize}
\item \texttt{world\_frame == odom\_frame}: publishes \texttt{odom}$\to$\texttt{base\_link}
\item \texttt{world\_frame == map\_frame}: publishes \texttt{map}$\to$\texttt{odom} (computed as $T_{\text{map}}^{\text{base}} \cdot (T_{\text{odom}}^{\text{base}})^{-1}$)
\end{itemize}

\textbf{Intuition:} \emph{This is the most important frame parameter.} Sets which role the EKF plays --- local smoother or global corrector.

\textbf{Effect:} The dual-EKF architecture exists exactly because of this parameter. Two filters, two values, decoupled transforms.

\textbf{Failure mode:} Both EKFs running with same \texttt{world\_frame} $\to$ TF conflict.
\end{paramcard}

% ============================================================
\section{Timing Parameters}
% ============================================================

\begin{paramcard}{frequency}
\textbf{Type:} double, Hz \quad \textbf{Project value:} 50.0 (EKFs) / 10.0 (navsat)

\textbf{What:} How often the EKF runs its predict step and publishes output.

\textbf{Math role:} Determines $\Delta t$ used in:
\begin{itemize}
\item Motion model integration: $\bar{\mu}_t = g(u_t, \mu_{t-1})$ uses $\Delta t = 1/\text{frequency}$
\item Process noise scaling: $\bar{\Sigma}_t = G_t\, \Sigma_{t-1}\, G_t^\top + R_t \cdot \Delta t$
\item Jacobian evaluation point and step size
\end{itemize}

\textbf{Intuition:} ``How often does the filter take a heartbeat?'' Higher = smoother output, more CPU. Lower = jerkier, lighter.

\textbf{Effect of increasing (e.g., 50 $\to$ 100 Hz):}
\begin{itemize}
\item Smaller $\Delta t$ per step $\to$ less linearisation error per step (Taylor expansion more accurate over smaller interval)
\item More frequent predictions absorb sensor jitter better
\item CPU load increases linearly
\item Diminishing returns past sensor rates --- predicting faster than sensors update gains nothing
\end{itemize}

\textbf{Effect of decreasing (e.g., 50 $\to$ 10 Hz):}
\begin{itemize}
\item Larger $\Delta t$ per step $\to$ more linearisation error
\item Output less smooth --- controller may oscillate
\item For nonlinear motion ($\omega \neq 0$), arc segments become noticeable
\end{itemize}

\textbf{Project rationale:} 50 Hz matches the IMU rate (typically 100 Hz raw $\to$ filtered to 50). Higher frequency than fastest sensor wastes CPU. \texttt{navsat\_transform} runs at 10 Hz to match the GPS update rate.

\textbf{Tuning rule:} Set to the rate of your fastest critical sensor, or what the controller expects (Nav2 typically wants $\geq 30$ Hz odometry).
\end{paramcard}

\begin{paramcard}{sensor\_timeout}
\textbf{Type:} double, seconds \quad \textbf{Project value:} 0.1

\textbf{What:} If no measurement arrives within this window, the EKF skips correction this cycle and just predicts.

\textbf{Math role:} Controls when the filter falls into \emph{prediction-only} mode. After timeout:
\[
\mu_t = \bar{\mu}_t = g(u_t, \mu_{t-1}), \quad \Sigma_t = \bar{\Sigma}_t = G_t \Sigma_{t-1} G_t^\top + R_t
\]
without any correction step. Uncertainty grows unboundedly during this period.

\textbf{Intuition:} ``How long do I wait for a sensor before giving up?'' Trade-off: too short = drops occasional late messages; too long = filter waits for stale data, output stale.

\textbf{Effect:}
\begin{itemize}
\item Short (0.05 s): aggressive, predicts often, robust to dropped messages
\item Long (0.5 s): conservative, may freeze on first sensor disconnect
\end{itemize}

\textbf{Project value 0.1 s:} matches 1/frequency = 1/50 = 0.02 s with 5x slack. Sensor offline for $>$ 0.1 s $\to$ predict-only.

\textbf{Failure mode:} Set too low with jittery sensors $\to$ filter stutters.
\end{paramcard}

\begin{paramcard}{transform\_time\_offset}
\textbf{Type:} double, seconds \quad \textbf{Project value:} 0.0

\textbf{What:} Adds a future timestamp offset to published TF, so other nodes can do future-time lookups.

\textbf{Math role:} Pure metadata --- shifts the timestamp on the published \texttt{TransformStamped}. State estimate unchanged.

\textbf{Intuition:} ``Backdate or future-date my TF transforms.'' Useful when consumers want to predict slightly into the future.

\textbf{Effect:}
\begin{itemize}
\item 0.0 (project): TF timestamp = filter timestamp. Standard.
\item 0.1: published TF claims to be valid for 0.1 s into future. Reduces ``extrapolation'' warnings in TF lookups by Nav2.
\end{itemize}
\end{paramcard}

\begin{paramcard}{transform\_timeout}
\textbf{Type:} double, seconds \quad \textbf{Project value:} 0.0

\textbf{What:} How long the TF buffer holds onto the previous transform if a new one is delayed.

\textbf{Math role:} TF interpolation/extrapolation buffer.

\textbf{Effect:} 0.0 means no timeout --- always serve latest. Tweak if downstream consumers are missing transforms.
\end{paramcard}

% ============================================================
\section{Mode Parameters}
% ============================================================

\begin{paramcard}{two\_d\_mode}
\textbf{Type:} bool \quad \textbf{Project value:} \texttt{true}

\textbf{What:} Constrains state to planar motion. $z, \phi, \theta, \dot{z}, \dot{\phi}, \dot{\theta}, \ddot{z}$ are clamped to zero.

\textbf{Math role:} Effectively reduces the 15-D state to an 8-D state $[x, y, \psi, \dot{x}, \dot{y}, \dot{\psi}, \ddot{x}, \ddot{y}]^\top$. Internally implemented by zero-ing rows/columns 2, 3, 4, 8, 9, 10, 14 of $\bar{\Sigma}_t$ after each update.

\textbf{Intuition:} ``This robot doesn't fly or fall.'' Skips degrees of freedom your robot physically can't have.

\textbf{Effect of \texttt{true}:}
\begin{itemize}
\item Prevents IMU pitch/roll noise from corrupting yaw estimate via cross-correlation
\item Prevents drift in $z$ from accumulated accelerometer error
\item Reduces effective state dimension $\to$ better-conditioned filter
\end{itemize}

\textbf{Effect of \texttt{false}:}
\begin{itemize}
\item Full 6-DOF tracking --- needed for drones, climbing robots, ramps
\item Roll/pitch from IMU integrated. Useful but introduces noise into $\psi$ via Coriolis/correlation
\end{itemize}

\textbf{Project rationale:} Lynx is a ground robot on flat-ish terrain. \texttt{true} is correct. If running on steep slopes or stairs, set to \texttt{false}.

\textbf{Failure mode:} Setting \texttt{true} when robot actually pitches significantly $\to$ accelerometer measurements get rotated incorrectly (gravity not removed properly), $\to$ horizontal acceleration biases.
\end{paramcard}

\begin{paramcard}{publish\_tf}
\textbf{Type:} bool \quad \textbf{Project value:} \texttt{true}

\textbf{What:} Whether the EKF publishes the world\_frame$\to$base\_link\_frame TF.

\textbf{Effect:} Both EKFs in this project publish their TF (one publishes \texttt{odom}$\to$\texttt{base\_link}, the other \texttt{map}$\to$\texttt{odom}). Setting \texttt{false} means another node must publish that TF (rare).

\textbf{Failure mode:} Two nodes both publishing same TF $\to$ TF tree warnings, intermittent values.
\end{paramcard}

\begin{paramcard}{publish\_acceleration}
\textbf{Type:} bool \quad \textbf{Project value:} \texttt{false}

\textbf{What:} Publishes filtered linear acceleration on \texttt{/accel/filtered}.

\textbf{Effect:} If \texttt{true}, additional message published with $\ddot{x}, \ddot{y}, \ddot{z}$ from state. If consumer cares (e.g., for vibration analysis), enable. Otherwise \texttt{false} saves bandwidth.
\end{paramcard}

\begin{paramcard}{print\_diagnostics}
\textbf{Type:} bool \quad \textbf{Project value:} \texttt{true}

\textbf{What:} Publishes diagnostic info on \texttt{/diagnostics} including measurement queue lengths, transform errors, etc.

\textbf{Effect:} Enabling helps tuning. No runtime cost worth worrying about.
\end{paramcard}

\begin{paramcard}{debug}
\textbf{Type:} bool \quad \textbf{Project value:} \texttt{false}

\textbf{What:} Verbose logging of every measurement, prediction, and correction.

\textbf{Effect:} \texttt{true} produces large debug files. Useful for tuning or post-mortem. Disable in production.
\end{paramcard}

% ============================================================
\section{Sensor Input Configuration}
% ============================================================

For each sensor (\texttt{odom0}, \texttt{odom1}, \texttt{imu0}, \ldots), several parameters control how its data enters the filter.

\subsection{Topic Name (\texttt{odomN}, \texttt{imuN}, \texttt{poseN}, \texttt{twistN})}

\begin{paramcard}{odom0 / imu0 / ...}
\textbf{Type:} string \quad \textbf{Project values:} \texttt{/diff\_drive\_controller/odom}, \texttt{/imu/data}, \texttt{/odometry/gps}, \texttt{/odometry/visual}, \texttt{/odometry/lidar}

\textbf{What:} ROS topic the EKF subscribes to.

\textbf{Math role:} Each topic provides one stream of measurements $z^{(s)}_t$ with covariance $Q^{(s)}_t$ (where $s$ indexes the sensor). The filter applies the standard EKF correction once per arriving message:
\[
K_t = \bar{\Sigma}_t H_t^\top (H_t \bar{\Sigma}_t H_t^\top + Q^{(s)}_t)^{-1}
\]

\textbf{Effect:} Each sensor enters as an independent correction step. Order of subscriber declaration in YAML doesn't matter; messages are processed in timestamp order.

\textbf{Best practice:} One sensor per topic. Don't merge two encoders into one topic --- their covariances become unrecoverable.
\end{paramcard}

\subsection{15-Bit Configuration Array}

This is the most important per-sensor parameter. It selects which state components the sensor measures.

\begin{paramcard}{odomN\_config / imuN\_config}
\textbf{Type:} list of 15 booleans \quad \textbf{Project examples below}

\textbf{What:} 15-element boolean array selecting which state dimensions this sensor contributes to. Mapping matches the state vector indexing.

\[
\texttt{config} = [b_x, b_y, b_z,\; b_\phi, b_\theta, b_\psi,\; b_{\dot{x}}, b_{\dot{y}}, b_{\dot{z}},\; b_{\dot{\phi}}, b_{\dot{\theta}}, b_{\dot{\psi}},\; b_{\ddot{x}}, b_{\ddot{y}}, b_{\ddot{z}}]
\]

\textbf{Math role:} Defines the observation matrix $H_t$ (or $C_t$ for linear case). Each \texttt{true} bit becomes a row of $H_t$ that picks out the corresponding state element. Example: if config selects only $\dot{x}$ and $\dot{\psi}$:
\[
H_t = \begin{bmatrix} 0 & 0 & 0 & 0 & 0 & 0 & \boxed{1} & 0 & 0 & 0 & 0 & 0 & 0 & 0 & 0 \\ 0 & 0 & 0 & 0 & 0 & 0 & 0 & 0 & 0 & 0 & 0 & \boxed{1} & 0 & 0 & 0 \end{bmatrix}
\]
The measurement equation becomes:
\[
z_t = \begin{bmatrix}\dot{x} \\ \dot{\psi}\end{bmatrix} + \delta_t
\]
The corresponding entries of the message's covariance matrix are extracted to form $Q_t$ for this update.

\textbf{Intuition:} ``Tell the filter what each sensor knows.'' A wheel encoder doesn't measure absolute position; only velocities. A GPS doesn't measure orientation. An IMU doesn't measure position.

\textbf{Project values walked through:}

\textbf{Wheel odometry (EKF \#1, \#2, \#3):}
\begin{lstlisting}[language=bash]
odom0_config: [false, false, false,    # x, y, z          NO
               false, false, false,    # roll, pitch, yaw NO
               true,  false, false,    # vx, vy, vz       vx YES
               false, false, true,     # vroll, vp, vyaw  vyaw YES
               false, false, false]    # ax, ay, az       NO
\end{lstlisting}
\textbf{Why:} Wheel encoders give body-frame linear velocity (forward only for diff-drive: $\dot{y}_{\text{body}} = 0$) and yaw rate. Position from \texttt{diff\_drive\_controller} drifts; we let the EKF integrate velocities itself, fusing with IMU.

\textbf{IMU (all EKFs):}
\begin{lstlisting}[language=bash]
imu0_config: [false, false, false,     # NO position
              true,  true,  true,      # roll, pitch, yaw YES
              false, false, false,     # NO linear vel
              true,  true,  true,      # ALL angular vels YES
              true,  true,  true]      # ALL linear accels YES
\end{lstlisting}
\textbf{Why:} IMU provides absolute roll/pitch from gravity (always reliable), yaw from gyro integration (drifts but high-rate), all angular velocities (raw gyro), and linear accelerations (with gravity removed). Note: yaw can be problematic if magnetometer is bad --- some setups disable bit \#5.

\textbf{GPS via navsat (EKF \#2):}
\begin{lstlisting}[language=bash]
odom1_config: [true,  true,  false,    # x, y YES (planar)
               false, false, false,
               false, false, false,
               false, false, false,
               false, false, false]
\end{lstlisting}
\textbf{Why:} GPS gives only absolute $(x, y)$ in map frame. No orientation, no velocity. $z$ disabled (\texttt{zero\_altitude} on navsat).

\textbf{Visual / LiDAR odometry (EKF \#3):}
\begin{lstlisting}[language=bash]
odom1_config: [true,  true,  false,    # x, y YES
               false, false, true,     # yaw YES
               false, false, false,
               false, false, false,
               false, false, false]
\end{lstlisting}
\textbf{Why:} Both VO and LO produce 6-DOF poses. In 2D mode we use $(x, y, \psi)$. Velocities not enabled to avoid double-counting (the wheel odom already provides those).

\textbf{Failure mode:} Enabling both pose and twist of the same sensor $\to$ filter receives a measurement and its time-derivative as ``independent,'' artificially over-confident.
\end{paramcard}

\subsection{Differential Mode}

\begin{paramcard}{odomN\_differential}
\textbf{Type:} bool \quad \textbf{Project value:} \texttt{false}

\textbf{What:} If \texttt{true}, the filter treats the difference between consecutive messages as a relative pose change rather than an absolute pose.

\textbf{Math role:} Instead of using $z_t = h(\bar{\mu}_t)$ as absolute pose, it uses:
\[
z_t \to (z_t - z_{t-1}) + \mu_{t-1}^{(\text{pose})}
\]
i.e., the new pose is the previous filter pose plus the measurement increment.

\textbf{Intuition:} ``Use this sensor's deltas, not its absolute readings.'' Useful when:
\begin{itemize}
\item Two sources both report absolute pose but disagree on origin (e.g., two SLAM systems)
\item You want to avoid GPS-induced jumps in the local filter
\end{itemize}

\textbf{Effect:}
\begin{itemize}
\item \texttt{false} (project): use absolute readings. Standard for one trusted source per modality.
\item \texttt{true}: differentiate first. Removes constant offsets. Doubles noise (subtraction increases variance).
\end{itemize}

\textbf{Project rationale:} \texttt{false} everywhere. Wheel odom feeds velocities (already differential by nature). GPS feeds map-frame position which is what we want. Visual/LiDAR feed pose deltas via separate node integration.

\textbf{Failure mode:} Setting \texttt{true} on velocity feeds re-differentiates them $\to$ acceleration measurements with huge variance.
\end{paramcard}

\begin{paramcard}{odomN\_relative}
\textbf{Type:} bool \quad \textbf{Project value:} \texttt{false}

\textbf{What:} If \texttt{true}, treats the first message as the origin --- subsequent messages are offsets from that first pose.

\textbf{Effect:} Useful when sensor reports in its own arbitrary frame. Project does not need this since wheel odom starts at origin and GPS starts at datum.
\end{paramcard}

\subsection{Queue and Latency}

\begin{paramcard}{odomN\_queue\_size}
\textbf{Type:} int \quad \textbf{Project value:} 10

\textbf{What:} ROS subscriber queue size for buffering incoming messages.

\textbf{Effect:} Larger absorbs bursts. If sensor publishes at 50 Hz and EKF processes at 50 Hz, queue of 10 holds 0.2 s buffer.

\textbf{Failure mode:} Queue too small + slow processing $\to$ messages dropped (visible in \texttt{/diagnostics}).
\end{paramcard}

\begin{paramcard}{odomN\_nodelay}
\textbf{Type:} bool \quad \textbf{Project value:} \texttt{true}/\texttt{false} (varies)

\textbf{What:} Sets TCP\_NODELAY on the subscription, disabling Nagle's algorithm to send small messages immediately.

\textbf{Effect:} \texttt{true} reduces latency for small high-rate streams (wheel odom). \texttt{false} default for IMU which already batches.
\end{paramcard}

% ============================================================
\section{Measurement Rejection Thresholds}
% ============================================================

These implement Mahalanobis gating --- rejecting measurements that are too far from the predicted mean.

\begin{paramcard}{imuN / odomN \_pose\_rejection\_threshold (and \_twist / \_linear\_acceleration variants)}
\textbf{Type:} double, Mahalanobis distance squared \quad \textbf{Project values:} 0.8 (IMU), 3.0 (GPS pose)

\textbf{What:} Maximum Mahalanobis distance for accepting a measurement. Larger distance $\to$ measurement rejected as outlier.

\textbf{Math role:} For each candidate measurement $z_t$, compute:
\[
d_M^2 = (z_t - h(\bar{\mu}_t))^\top S_t^{-1} (z_t - h(\bar{\mu}_t)), \qquad S_t = H_t \bar{\Sigma}_t H_t^\top + Q_t
\]
Reject if $d_M^2 > \text{threshold}$.

\textbf{Intuition:} ``How surprised do I have to be before I assume the sensor is lying?'' Smaller = stricter, more rejections. Larger = more tolerant, accepts noisy data.

\textbf{Statistical interpretation:} Under the assumption that innovation is zero-mean Gaussian with covariance $S_t$, $d_M^2$ follows $\chi^2_k$ where $k$ is measurement dimension. Common thresholds:
\begin{center}
\begin{tabular}{lll}
\toprule
$k$ & 95\% threshold & 99\% threshold \\
\midrule
1 & 3.84 & 6.63 \\
2 & 5.99 & 9.21 \\
3 & 7.81 & 11.34 \\
6 & 12.59 & 16.81 \\
\bottomrule
\end{tabular}
\end{center}

\textbf{Project values explained:}
\begin{itemize}
\item IMU 0.8: very strict. Reflects high confidence in IMU --- big residual = something else wrong.
\item GPS 3.0: moderate. GPS is noisy by nature; allow some surprise before rejecting.
\end{itemize}

\textbf{Effect of decreasing (e.g., 3.0 $\to$ 1.0):}
\begin{itemize}
\item More GPS measurements rejected as outliers
\item Faster recovery from rare GPS spikes (multipath, etc.)
\item BUT: filter under-corrects when state has genuinely diverged --- can lock to drifted estimate
\end{itemize}

\textbf{Effect of increasing (e.g., 3.0 $\to$ 100.0):}
\begin{itemize}
\item All measurements accepted
\item Outliers contaminate state
\item Useful initially during tuning; tighten later
\end{itemize}

\textbf{Failure mode:} Setting too low + initial covariance set very small (overconfident initial state) $\to$ all measurements rejected, filter drifts forever.
\end{paramcard}

% ============================================================
\section{Process Noise Covariance $R_t$}
% ============================================================

This is the most-tuned parameter. It controls how much the filter ``trusts itself.''

\begin{paramcard}{process\_noise\_covariance}
\textbf{Type:} 15$\times$15 matrix (225 floats), specified row-major \quad \textbf{Project values: see below}

\textbf{What:} Covariance matrix $R_t$ added to predicted covariance each step. Represents unmodelled state changes per unit time.

\textbf{Math role:} In the predict step:
\[
\bar{\Sigma}_t = G_t\, \Sigma_{t-1}\, G_t^\top + R_t \cdot \Delta t
\]
(robot\_localization scales $R$ by $\Delta t$ internally so values are per-second.)

\textbf{Intuition:} ``How much can the world change between predictions, beyond what my motion model captures?'' Models:
\begin{itemize}
\item Slip and unmodelled dynamics
\item Wheel radius variation
\item Calibration error
\item External forces (wind, ground irregularities)
\end{itemize}

\textbf{Effect of increasing (e.g., $R_{xx}: 0.05 \to 0.5$):}
\begin{itemize}
\item Predicted covariance $\bar{\Sigma}_t$ grows faster
\item Kalman gain $K_t = \bar{\Sigma}_t H_t^\top S_t^{-1}$ grows
\item Filter trusts measurements more, follows them faster
\item Output noisier (less smoothing)
\end{itemize}

\textbf{Effect of decreasing (e.g., $R_{xx}: 0.05 \to 0.005$):}
\begin{itemize}
\item Predicted covariance grows slowly
\item Kalman gain small --- filter ignores measurements
\item Output very smooth but slow to respond
\item Risk: filter believes itself even when wrong (divergence)
\end{itemize}

\textbf{Project EKF \#1 (odom) diagonal --- exact values:}
\begin{center}
\footnotesize
\begin{tabular}{cccc}
\toprule
Index & Component & $R_{ii}$ & Interpretation \\
\midrule
0 & $x$ & 0.05 & 22 cm $/\sqrt{\text{s}}$ position uncertainty growth \\
1 & $y$ & 0.05 & same \\
2 & $z$ & 0.06 & higher (vertical bouncing on terrain) --- ignored in 2D mode \\
3 & $\phi$ (roll) & 0.03 & 0.17 rad $/\sqrt{\text{s}}$ \\
4 & $\theta$ (pitch) & 0.03 & same \\
5 & $\psi$ (yaw) & 0.06 & 0.24 rad $/\sqrt{\text{s}}$ --- larger because wheel-baseline error \\
6 & $\dot{x}$ & 0.025 & 0.16 (m/s)/$\sqrt{\text{s}}$ --- forward vel uncertainty \\
7 & $\dot{y}$ & 0.025 & holonomic constraint enforces, but kept symmetric \\
8 & $\dot{z}$ & 0.04 & vertical (ignored in 2D) \\
9 & $\dot{\phi}$ & 0.01 & angular vel small uncertainty \\
10 & $\dot{\theta}$ & 0.01 & same \\
11 & $\dot{\psi}$ & 0.02 & yaw rate uncertainty \\
12 & $\ddot{x}$ & 0.01 & accel uncertainty \\
13 & $\ddot{y}$ & 0.01 & same \\
14 & $\ddot{z}$ & 0.015 & gravity residual \\
\bottomrule
\end{tabular}
\end{center}

\textbf{Project EKF \#2 (map) diagonal --- post-debug final values:}
\begin{center}
\footnotesize
\begin{tabular}{ccccp{4.7cm}}
\toprule
Index & Component & $R_{ii}$ & vs EKF\#1 & Reason \\
\midrule
0  & $x$           & 0.05  & same  & RTK GPS at 1cm $\sigma$ pulls fast; moderate $Q$ stops state lock-in \\
1  & $y$           & 0.05  & same  & same \\
5  & $\psi$ (yaw)  & 0.001 & \textbf{60$\times$ tighter} & lock yaw, kill $P[\psi,x]$ coupling that drags yaw on GPS x,y updates \\
6  & $\dot{x}$     & 0.001 & \textbf{25$\times$ tighter} & non-holonomic: $v_x \approx 0$ unless commanded \\
7  & $\dot{y}$     & 0.001 & \textbf{25$\times$ tighter} & non-holonomic: $v_y \equiv 0$ for diff-drive \\
11 & $\dot{\psi}$  & 0.001 & \textbf{20$\times$ tighter} & lock $\omega_z$, anchored by wheel + LiDAR $\Delta\theta$ \\
\bottomrule
\end{tabular}
\end{center}

\textbf{Why EKF \#2 needs tight velocity / yaw $R$ (post-debug rationale):}
\begin{itemize}
\item Old reasoning (loose $R$) assumed consumer GPS at $\sigma\!\approx\!0.5$~m. With RTK at $\sigma\!=\!1$~cm, filter no longer needs to ``flex'' to absorb large GPS jumps.
\item Loose $Q[\psi]$ lets GPS x,y updates pull yaw via cross-covariance $P[\psi, x]$. Measured: 22$^\circ$ stationary yaw walk in 50~s with $Q[\psi]=0.06$.
\item Loose $Q[\dot{y}]$ allows $v_y$ state to drift; on diff-drive that is unphysical. Measured: $v_y$ ran to $+7.33$~m/s before being clamped by GPS feedback. Caused 18~m Y drift.
\item GPS rejection threshold raised to 100$\sigma$ to ensure RTK is always trusted (no realistic outlier scenario at 1~cm noise).
\end{itemize}

\textbf{Wheel $v_y$ row enabled.} \texttt{odom0\_config[7]: false $\to$ true}. \texttt{diff\_drive\_controller} publishes $v_y\!=\!0$ with cov $0.001$. Without fusing this, EKF~\#2 had no measurement source for $v_y$; only IMU accel (now disabled) or process noise constrained it. Result was velocity runaway.

\textbf{IMU linear acceleration disabled in EKF~\#2.} \texttt{imu0\_config[12..14]: false}. Even with bias\_mean=0, residual accel noise integrates to velocity drift. Wheel + GPS provide all needed constraints in map frame; accel adds no information.

\textbf{Tuning rule of thumb:}
\begin{enumerate}
\item Start with project values
\item Drive a known trajectory
\item Plot innovation: $\nu_t = z_t - h(\bar{\mu}_t)$
\item If innovation grows over time (filter drifts) $\to$ increase $R$
\item If innovation noisy/jumps (filter overshooting) $\to$ decrease $R$ or increase $Q$
\item Validate with NIS test (see Validation section)
\end{enumerate}

\textbf{Off-diagonal entries:} Project keeps $R_t$ diagonal (zero off-diagonals). This assumes process noise is uncorrelated across state dimensions. A more accurate model could have $R_{x \dot{x}} \neq 0$ since position and velocity errors are coupled. \texttt{robot\_localization} accepts this but most implementations leave it diagonal.
\end{paramcard}

\begin{intuitionbox}{Why Yaw $R_{\psi\psi}$ is 0.06 (Twice $R_{\phi\phi}$)}
The wheel-encoder yaw estimate uses $\omega = (v_R - v_L) / b$ where $b$ is the baseline. A 1\% baseline error becomes a 1\% systematic yaw rate error --- accumulating linearly over distance. Roll and pitch from IMU gravity vector have no such systematic source --- only zero-mean noise. Hence the higher process noise on yaw acknowledges that even with a perfect IMU, yaw can drift due to wheel calibration error.
\end{intuitionbox}

% ============================================================
\section{Initial Estimate Covariance $\Sigma_0$}
% ============================================================

\begin{paramcard}{initial\_estimate\_covariance}
\textbf{Type:} 15$\times$15 matrix \quad \textbf{Project values: see below}

\textbf{What:} The covariance $\Sigma_0$ at filter startup, before any measurement.

\textbf{Math role:} Just the initial value of $\Sigma_t$ for $t = 0$. Affects the early Kalman gain:
\[
K_0 = \Sigma_0 H_0^\top (H_0 \Sigma_0 H_0^\top + Q_0)^{-1}
\]
Large $\Sigma_0 \to$ filter trusts first measurement strongly. Small $\Sigma_0 \to$ filter resists initial measurements.

\textbf{Intuition:} ``How much do I know about my starting state?'' If you start at known origin, $\Sigma_0 \approx 0$ (use \texttt{1e-9}). If you start at unknown location, $\Sigma_0$ large.

\textbf{Project EKF \#1 (odom):} all diagonal entries \texttt{1e-9}. \\
\textbf{Why:} The robot is assumed to start at \texttt{(0, 0, 0)} in odom frame. Origin is the robot's startup pose by definition.

\textbf{Project EKF \#2 (map):} diagonal \texttt{[1.0, 1.0, 1e-9, 1.0, 1.0, 1e-9, 1.0, 1.0, ..., 1.0]} \\
\textbf{Why:} Map-frame position $(x, y)$ initially \emph{unknown} until first GPS fix. $\Sigma_{xx}^{(0)} = 1.0$ m$^2$ allows the first GPS measurement (covariance 0.26 m$^2$) to dominate via Kalman gain:
\[
K \approx \frac{1.0}{1.0 + 0.26} \approx 0.79
\]
Filter snaps to GPS quickly. After a few fixes, $\Sigma$ shrinks below GPS variance; gain drops.

\textbf{Effect of decreasing initial covariance below sensor noise:} filter ``trusts'' an unknown initial pose more than the actual measurement --- output anchored to wrong initial value, drifts away.

\textbf{Effect of increasing way past needed:} no harm; filter converges to correct value within a few cycles regardless.

\textbf{Tuning rule:} Set diagonal entries to \texttt{1e-9} for components you start at zero (or known value), and to a few $\sigma^2$ of the sensor that bounds them otherwise.
\end{paramcard}

% ============================================================
\section{Control Input Parameters (Local EKF Only)}
% ============================================================

EKF \#1 and \#3 use \texttt{cmd\_vel} as a control input to improve prediction between sensor updates.

\begin{paramcard}{use\_control}
\textbf{Type:} bool \quad \textbf{Project value:} \texttt{true} (EKF \#1, \#3) / \texttt{false} (EKF \#2)

\textbf{What:} Whether the EKF subscribes to \texttt{cmd\_vel} and uses it as control input $u_t$ in the motion model.

\textbf{Math role:} With control, the prediction is:
\[
\bar{\mu}_t = g(u_t, \mu_{t-1})
\]
where $g$ depends on commanded $(v, \omega)$. Without control, $u_t$ is treated as zero or constant velocity --- prediction degenerates to constant-velocity model.

\textbf{Intuition:} ``I know what I told the robot to do; use that to predict.'' Helps when sensors are silent or stale.

\textbf{Effect:}
\begin{itemize}
\item \texttt{true}: prediction follows command, filter snappier
\item \texttt{false}: relies entirely on past state + sensors. Smoother in steady-state but worse during transients.
\end{itemize}

\textbf{Project rationale:} Local EKF benefits from \texttt{cmd\_vel} because wheel feedback has latency. Global EKF doesn't because it operates on already-fused local odometry.
\end{paramcard}

\begin{paramcard}{control\_config}
\textbf{Type:} list of 6 booleans, $[\dot{x}, \dot{y}, \dot{z}, \dot{\phi}, \dot{\theta}, \dot{\psi}]$ \quad \textbf{Project value:} \texttt{[true, false, false, false, false, true]}

\textbf{What:} Which fields of \texttt{cmd\_vel} are valid commands.

\textbf{Project value:} only $\dot{x}$ (forward velocity) and $\dot{\psi}$ (yaw rate). Diff-drive constraint --- robot cannot command lateral velocity.
\end{paramcard}

\begin{paramcard}{acceleration\_limits / deceleration\_limits}
\textbf{Type:} list of 6 doubles, m/s$^2$ or rad/s$^2$ \quad \textbf{Project value:} \texttt{[1.0, 0, 0, 0, 0, 1.5]}

\textbf{What:} Maximum (de)acceleration the robot can physically achieve.

\textbf{Math role:} Used to limit how fast the predicted velocity can change toward the commanded velocity:
\[
\bar{v}_t = \mathrm{clamp}(\bar{v}_t, v_{t-1} - a_{\max}\Delta t, v_{t-1} + a_{\max}\Delta t)
\]

\textbf{Intuition:} ``This robot has mass. It can't change velocity instantly.''

\textbf{Project values:} 1.0 m/s$^2$ linear, 1.5 rad/s$^2$ angular --- modest values matching small ground robot.

\textbf{Failure mode:} Set too low $\to$ filter doesn't predict sharp commanded changes, lags. Set too high $\to$ filter believes robot can teleport, doesn't smooth.
\end{paramcard}

\begin{paramcard}{acceleration\_gains / deceleration\_gains}
\textbf{Type:} list of 6 doubles in $[0, 1]$ \quad \textbf{Project value:} \texttt{[0.8, 0, 0, 0, 0, 0.9]}

\textbf{What:} How fast the robot is expected to ramp acceleration toward commanded value. 1.0 = instant, 0.0 = no response.

\textbf{Math role:} Models actuator dynamics. Predicted acceleration:
\[
a_t = g_a \cdot a_{\text{commanded}}
\]

\textbf{Intuition:} ``Even commanded acceleration is not delivered fully --- there's a control loop in the motors.''

\textbf{Project values:} 0.8 linear, 0.9 yaw --- 80--90\% of commanded acceleration is realised.
\end{paramcard}

\begin{paramcard}{control\_timeout}
\textbf{Type:} double, seconds \quad \textbf{Project value:} 0.2

\textbf{What:} If no \texttt{cmd\_vel} arrives within this time, control input is treated as zero.

\textbf{Effect:} Prevents stale commands from biasing prediction when teleop disconnects.
\end{paramcard}

% ============================================================
\section{Why GPS is Excluded from EKF \#2 (Covariance Cross-Term Leak)}
\label{sec:gps-exclusion}
% ============================================================

This section explains the engineering reason behind this project's most counter-intuitive design decision: removing GPS from EKF \#2. Read carefully before re-enabling.

\subsection{The Cross-Covariance Problem}

The EKF state covariance $\Sigma$ is a dense $15 \times 15$ matrix. Off-diagonal entries $\Sigma_{x,\psi}$ and $\Sigma_{y,\psi}$ encode statistical correlation between position and yaw --- ``if I learn something about $x$, what does it tell me about $\psi$?''

These cross-terms become non-zero whenever the robot moves with $\omega \neq 0$. The motion-model Jacobian:
\[
G_t = \begin{bmatrix} 1 & 0 & -v\sin(\mu_\theta)\Delta t \\ 0 & 1 & v\cos(\mu_\theta)\Delta t \\ 0 & 0 & 1 \end{bmatrix}
\]
has non-zero $\partial g_x / \partial \psi$ and $\partial g_y / \partial \psi$. Predicted covariance $\bar{\Sigma}_t = G_t\,\Sigma_{t-1}\,G_t^\top + R_t$ inherits these cross-terms. After any curved motion, $\Sigma_{x,\psi}, \Sigma_{y,\psi} \neq 0$.

\subsection{The Leak Path}

GPS measures only $[x, y]$. Observation Jacobian:
\[
H_{\text{GPS}} = \begin{bmatrix} 1 & 0 & 0 & 0 & 0 & 0 & \cdots \\ 0 & 1 & 0 & 0 & 0 & 0 & \cdots \end{bmatrix}
\]

Kalman gain $K_t = \bar{\Sigma}_t H_t^\top S_t^{-1}$ has a row for \emph{every} state component, including yaw. The yaw-row entries:
\[
K_{\psi, x} = \frac{\Sigma_{\psi,x}}{\Sigma_{xx} + Q_{\text{GPS}}}, \qquad K_{\psi, y} = \frac{\Sigma_{\psi,y}}{\Sigma_{yy} + Q_{\text{GPS}}}
\]

Yaw correction:
\[
\Delta\psi = K_{\psi,x}\,\nu_x + K_{\psi,y}\,\nu_y
\]
where $\nu = z_{\text{GPS}} - h(\bar{\mu}_t)$ is the innovation.

\textbf{GPS noise of 0.5 m produces non-zero innovation every cycle. Yaw moves --- even though GPS never observed yaw.}

\subsection{Numerical Example (Project Values)}

Realistic state after some driving:
\[
\Sigma_{xx} = 0.30 \text{ m}^2, \quad \Sigma_{\psi\psi} = 0.02 \text{ rad}^2, \quad \Sigma_{x,\psi} = 0.04
\]

GPS noise spike $\nu_x = 0.5$ m, $Q_{\text{GPS}} = 0.26$ m$^2$ (your measured GPS variance):
\[
K_{\psi,x} = \frac{0.04}{0.30 + 0.26} \approx 0.071
\]
\[
\Delta\psi = 0.071 \cdot 0.5 \approx 0.036 \text{ rad} \approx 2°
\]

A single GPS fix moved yaw by $2°$. At 10 Hz GPS, sustained walk emerges. Measured drift in this project: $3.56°/$s.

\subsection{Compounding via VO/LiDAR}

Joseph form covariance update:
\[
\Sigma_{\psi,x}^{\text{new}} = \Sigma_{\psi,x} - K_{\psi,x}\,\Sigma_{xx}
\]
reduces but doesn't zero the cross-term. Next predict step regrows it through $G_t$.

Worse: VO and LiDAR observe \emph{both} position and yaw. Their $H$ matrix has rows for $x, y, \psi$. The Kalman gain spreads VO/LiDAR's $x$-correction into yaw via the now-shaped covariance. Each multi-sensor cycle amplifies the leak.

\subsection{Why the Canonical Tutorial Doesn't Suffer}

Tom Moore's reference example assumes an absolute yaw observation:
\begin{itemize}
\item Dual-antenna RTK GPS providing measured heading
\item Reliable magnetometer with hard/soft-iron calibration
\item External heading reference (compass, sun sensor)
\end{itemize}

That observation has small $Q_\psi$. Kalman gain heavily weights the direct yaw measurement, dominating the indirect leak from cross-terms. Yaw stays bounded.

\textbf{This project has none of those.} Gazebo IMU has no magnetometer; ``absolute'' yaw is integrated gyro --- drifts $> 150°$ in 2 minutes by itself. Nothing fights the leak.

\subsection{The Sensor Consensus Solution}

Remove GPS from EKF \#2. Both EKFs now process identical sensor sets:
\[
\mu_t^{(\text{global})} \approx \mu_t^{(\text{local})} \implies T_{\text{map}}^{\text{odom}} \approx I
\]

Map frame stays near-identity to odom frame. No GPS-driven cross-term injection. Map frame angularly stable.

\textbf{Empirical evidence in this project:}
\begin{itemize}
\item Earlier config (GPS in EKF \#2): map yaw drifted at $3.56°/$s, accumulating $> 200°$ in 2 minutes --- catastrophic for navigation.
\item Current config (GPS removed): map$\to$odom transform stays within $\pm 0.1$ m and $\pm 1°$ over 2-minute run.
\item Disable GPS $\to$ drift disappears. Re-enable $\to$ drift returns.
\end{itemize}

\subsection{Trade-Off}

\begin{center}
\footnotesize
\begin{tabular}{p{4cm} p{4cm} p{4cm}}
\toprule
\textbf{Aspect} & \textbf{GPS in EKF \#2} & \textbf{GPS excluded (project)} \\
\midrule
Yaw stability & drifts $\sim 3.56°/$s & rock solid \\
Global anchoring & yes (cm-m accuracy) & no (drifts with odom) \\
Re-anchorable to lat/lon & yes & no, only spawn-relative \\
Best for & open outdoor + RTK + magnetometer & indoor / GPS-denied / short-range \\
\bottomrule
\end{tabular}
\end{center}

For most navigation tasks, a stable heading matters more than centimetric global position. A robot drifting $3.56°/$s of yaw will spin in circles trying to reach waypoints. A robot drifting 0.5 m/min of global position can still navigate a 50 m mission.

\subsection{If You Need Global Anchoring Back}

Four options:

\begin{enumerate}
\item \textbf{Add a strong yaw observation}: dual-antenna RTK, calibrated magnetometer (with hard/soft-iron calibration), or sun sensor. Strong $Q_\psi$ measurement drowns out the leak.

\item \textbf{Diagonalise $\Sigma$ post-GPS-update}: zero out $\Sigma_{x,\psi}$ and $\Sigma_{y,\psi}$ after each GPS correction. Hacky but effective. Requires custom EKF code (not available in stock \texttt{robot\_localization}).

\item \textbf{Loosely-coupled GPS}: run a separate filter that consumes only EKF \#2 output and GPS, publish a \texttt{utm}$\to$\texttt{map} transform without polluting EKF \#2's state.

\item \textbf{Periodic re-anchoring}: when GPS confidence is high (RTK fix), trigger a one-shot pose reset of EKF \#2 instead of continuous fusion. Resets cross-term to zero, breaks the compounding cycle.
\end{enumerate}

\subsection{Where to Look in YAML}

Compare the actual project files:
\begin{itemize}
\item \texttt{src/husarion\_ugv\_description/config/dual\_ekf\_navsat.yaml} --- canonical-style file with three EKFs, originally had GPS in EKF \#2
\item \texttt{src/husarion\_ugv\_description/config/baseline\_ekf.yaml} --- current production config, GPS not in EKF \#2 input list
\item \texttt{src/husarion\_ugv\_description/config/hybrid\_qf\_ekf.yaml} --- experimental quantum filter variant
\end{itemize}

The relevant block in EKF \#2's section (post-debug, final working):
\begin{lstlisting}[language=bash]
ekf_filter_node_map:
  ros__parameters:
    world_frame: map
    use_control: false

    # Wheel: vx, vy (kinematic constraint), vyaw
    odom0: /diff_drive_controller/odom
    odom0_config: [false, false, false,
                   false, false, false,
                   true,  true,  false,   # vx, vy, vz
                   false, false, true,    # vroll, vpitch, vyaw
                   false, false, false]
    odom0_differential: false

    # GPS: x, y absolute (RTK, ~1cm noise)
    odom1: /odometry/gps
    odom1_config: [true, true, false,
                   false, false, false,
                   false, false, false,
                   false, false, false,
                   false, false, false]
    odom1_pose_rejection_threshold: 100.0  # RTK can't outlier
    odom1_differential: false

    # LiDAR: yaw absolute, differential mode (Δθ between samples)
    odom2: /odometry/lidar
    odom2_config: [false, false, false,
                   false, false, true,    # yaw only
                   false, false, false,
                   false, false, false,
                   false, false, false]
    odom2_differential: true

    # IMU: angular velocity only (no accel - bias integrates to drift)
    imu0: /imu/data/covariance_fixed
    imu0_config: [false, false, false,
                  false, false, false,
                  false, false, false,
                  true,  true,  true,    # vroll, vpitch, vyaw
                  false, false, false]   # accel DISABLED
    imu0_differential: false

    # Tight process noise: lock yaw + velocities (non-holonomic)
    # Q[x,y]    = 0.05
    # Q[yaw]    = 0.001
    # Q[vx,vy]  = 0.001
    # Q[vyaw]   = 0.001
\end{lstlisting}

Launch file remap (the topology fix):
\begin{lstlisting}[language=bash]
# baseline_ekf.launch.py, navsat_transform node:
remappings=[
    ("imu",               "imu/data/covariance_fixed"),
    ("gps/fix",           "gps/fix/covariance_fixed"),
    ("odometry/filtered", "odometry/global"),  # NOT odometry/local
],
\end{lstlisting}

% ============================================================
\section{NavSat Transform Parameters}
% ============================================================

\begin{paramcard}{frequency (navsat)}
\textbf{Type:} double, Hz \quad \textbf{Project value:} 10.0

\textbf{What:} How often \texttt{navsat\_transform\_node} publishes \texttt{/odometry/gps}.

\textbf{Effect:} Should match GPS update rate. Higher wastes CPU; lower drops fixes.
\end{paramcard}

\begin{paramcard}{delay}
\textbf{Type:} double, seconds \quad \textbf{Project value:} 3.0

\textbf{What:} Wait this long after first GPS fix before computing the datum transform.

\textbf{Math role:} Allows the EKF to stabilise its initial estimate before locking the GPS-to-map alignment. The datum is computed from:
\[
T_{\text{map}}^{\text{UTM}} = \text{Yaw}(\psi_0) \cdot \text{Translate}(-x_0^{\text{UTM}}, -y_0^{\text{UTM}})
\]
where $\psi_0$ is taken from the IMU after \texttt{delay} seconds.

\textbf{Intuition:} ``Wait for everything to settle before snapshotting the world.''

\textbf{Effect of small delay:} initial heading from IMU may be noisy --- map frame misaligned.

\textbf{Effect of large delay:} robot may have moved before datum set --- origin offset.
\end{paramcard}

\begin{paramcard}{magnetic\_declination\_radians}
\textbf{Type:} double \quad \textbf{Project value:} 0.0

\textbf{What:} Angle between true north and magnetic north at robot's location.

\textbf{Math role:} Used to align the IMU's magnetometer-derived yaw with true north (UTM grid):
\[
\psi_{\text{true}} = \psi_{\text{magnetic}} + \delta_{\text{declination}}
\]

\textbf{Intuition:} ``Compass doesn't point to true north.'' Varies geographically: $-20^\circ$ to $+20^\circ$ globally.

\textbf{Project value 0.0:} acceptable in simulation (Gazebo provides true north). For real robot at Edinburgh: declination $\approx -1.2^\circ \approx -0.021$ rad. Look up at \url{https://www.ngdc.noaa.gov/geomag-web/}.

\textbf{Failure mode:} Wrong declination $\to$ map frame rotated $\to$ all GPS-derived positions twisted by that angle. Robot ``walks sideways'' relative to map.
\end{paramcard}

\begin{paramcard}{yaw\_offset}
\textbf{Type:} double, radians \quad \textbf{Project value:} 0.0

\textbf{What:} Rotation between IMU's zero-yaw direction and the conventional $+x$ axis (East in ENU).

\textbf{Math role:} If IMU's $\psi=0$ points North, but UTM East is the $+x$ axis, you need $\texttt{yaw\_offset} = -\pi/2$.

\textbf{Project value 0.0:} simulator IMU is configured so $\psi = 0$ aligns with $+x$ East.

\textbf{Failure mode:} Most common config error. Robot drives east according to filter but actually drives north. Diagnose by driving in commanded direction and checking GPS.
\end{paramcard}

\begin{paramcard}{zero\_altitude}
\textbf{Type:} bool \quad \textbf{Project value:} \texttt{true}

\textbf{What:} Force altitude in published odometry to zero, regardless of GPS altitude.

\textbf{Effect:} For a planar robot, this prevents GPS altitude noise (1.23 m$^2$ in this project!) from corrupting the $z$ estimate.

\textbf{Project rationale:} Lynx is a ground robot. Altitude irrelevant.
\end{paramcard}

\begin{paramcard}{publish\_filtered\_gps}
\textbf{Type:} bool \quad \textbf{Project value:} \texttt{true}

\textbf{What:} Publishes the EKF's filtered estimate back as a \texttt{NavSatFix} (lat/lon).

\textbf{Effect:} Useful for visualisation in mapping tools (Google Maps, Foxglove). No effect on filter behavior.
\end{paramcard}

\begin{paramcard}{use\_odometry\_yaw}
\textbf{Type:} bool \quad \textbf{Project value:} \texttt{false}

\textbf{What:} If \texttt{true}, takes initial heading from \texttt{/odometry/filtered} (i.e., wheel + IMU fused yaw) instead of raw IMU.

\textbf{Effect:}
\begin{itemize}
\item \texttt{false} (project): use IMU yaw at startup --- absolute reference (assuming magnetometer)
\item \texttt{true}: use fused yaw, which is initially zero (odom frame origin)
\end{itemize}

Project chose \texttt{false}: relies on IMU's absolute heading sensor for global alignment.
\end{paramcard}

\begin{paramcard}{wait\_for\_datum}
\textbf{Type:} bool \quad \textbf{Project value:} \texttt{false}

\textbf{What:} If \texttt{true}, waits for an explicit \texttt{datum} parameter before computing transforms. If \texttt{false}, uses first GPS fix as datum automatically.

\textbf{Effect:}
\begin{itemize}
\item \texttt{false} (project): convenient for tests --- start anywhere, first fix becomes origin
\item \texttt{true}: requires \texttt{datum: [lat, lon, yaw]} explicitly. Repeatable across runs --- map frame consistent across launches.
\end{itemize}

\textbf{Production setup:} should use \texttt{true} with a known datum so map frame is consistent across robot restarts.
\end{paramcard}

\begin{paramcard}{broadcast\_cartesian\_transform}
\textbf{Type:} bool \quad \textbf{Project value:} \texttt{true}

\textbf{What:} Whether to publish \texttt{utm}$\to$\texttt{world\_frame} TF.

\textbf{Effect:} Useful for tools that need UTM coordinates directly. Doesn't affect EKF.
\end{paramcard}

% ============================================================
\section{Per-Sensor Covariance Override}
% ============================================================

\begin{paramcard}{odomN\_pose\_covariance\_diagonal / odomN\_twist\_covariance\_diagonal}
\textbf{Type:} list of 6 doubles \quad \textbf{Project values: see below}

\textbf{What:} Override the message's reported covariance with these values.

\textbf{Math role:} Replaces the diagonal entries of $Q_t$ used in the update step:
\[
Q_t^{(s)} = \mathrm{diag}(\text{covariance\_diagonal})
\]

\textbf{Intuition:} ``Sensor's own covariance is wrong --- use my measured value instead.''

\textbf{Project values (EKF \#3 fusion):}
\begin{lstlisting}[language=bash]
# Wheel odom: tight pose/twist (very confident on velocities)
odom0_pose_covariance_diagonal:  [0.001, 0.001, 0.001, 0.001, 0.001, 0.001]
odom0_twist_covariance_diagonal: [0.001, 0.001, 0.001, 0.001, 0.001, 0.001]

# Visual odom: looser pose, looser velocity
odom1_pose_covariance_diagonal:  [0.01, 0.01, 0.01, 0.01, 0.01, 0.05]
odom1_twist_covariance_diagonal: [0.1, 0.1, 0.1, 0.1, 0.1, 0.1]

# LiDAR odom: tighter than VO
odom2_pose_covariance_diagonal:  [0.005, 0.005, 0.005, 0.01, 0.01, 0.02]
odom2_twist_covariance_diagonal: [0.1, 0.1, 0.1, 0.1, 0.1, 0.1]
\end{lstlisting}

\textbf{Why these values:}
\begin{itemize}
\item \textbf{Wheel} 0.001: trusted velocity source, very tight. Note: trusting position to 0.001 is meaningless because we disable position bits.
\item \textbf{Visual} 0.01--0.05: VO drifts with distance, monocular has scale issues. Larger yaw covariance (0.05) reflects rotation ambiguity.
\item \textbf{LiDAR} 0.005--0.02: more accurate than VO in structured environments. Tightest non-wheel source.
\end{itemize}

\textbf{Tuning order:} Use override only when you know the message's covariance is wrong (very common --- many drivers publish placeholder values like all-1s or all-0s).
\end{paramcard}

% ============================================================
\section{IMU-Specific Parameters}
% ============================================================

\begin{paramcard}{imuN\_remove\_gravitational\_acceleration}
\textbf{Type:} bool \quad \textbf{Project value:} \texttt{true}

\textbf{What:} Subtracts gravity from accelerometer readings before fusing.

\textbf{Math role:} Raw accel reads:
\[
\tilde{a}^b = a^b_{\text{true}} - R^b_w\, g
\]
With this flag \texttt{true}, the filter computes:
\[
a^b_{\text{linear}} = \tilde{a}^b + R^b_w\, g
\]
using the current orientation estimate to rotate gravity into body frame.

\textbf{Intuition:} ``Don't let gravity look like horizontal acceleration when the robot tilts.''

\textbf{Effect of \texttt{true}:}
\begin{itemize}
\item Filter receives only motion-induced acceleration
\item Robust to terrain pitch/roll
\end{itemize}

\textbf{Effect of \texttt{false}:}
\begin{itemize}
\item Filter receives raw $\tilde{a}^b$ including $-R^b_w\, g$
\item In 2D mode, planar accel includes gravity components when robot pitches
\item Sensor model $h$ would need to predict gravity itself
\end{itemize}

\textbf{Failure mode:} \texttt{false} on a sloped robot $\to$ filter believes constant horizontal acceleration $\to$ velocity drifts.
\end{paramcard}

% ============================================================
\section{The Big Picture: How Parameters Compose}
% ============================================================

\subsection{Causal Chain --- One Update Cycle}

Imagine GPS has just arrived. Trace through what happens:

\begin{enumerate}
\item \textbf{Predict step uses}:
\begin{itemize}
\item \texttt{frequency} sets $\Delta t$
\item \texttt{two\_d\_mode} clamps non-planar dims
\item \texttt{use\_control}, \texttt{control\_config}, \texttt{acceleration\_limits} shape $g(u_t, \mu_{t-1})$
\item \texttt{process\_noise\_covariance} provides $R_t$
\item Result: $\bar{\mu}_t, \bar{\Sigma}_t$
\end{itemize}

\item \textbf{Sensor message arrives on \texttt{/odometry/gps}}:
\begin{itemize}
\item \texttt{odom1} topic identifies it
\item \texttt{odom1\_config} extracts which dims to use $\to$ builds $H_t$
\item Message's pose covariance becomes $Q_t$ (or override via \texttt{odom1\_pose\_covariance\_diagonal})
\item \texttt{odom1\_differential, \_relative} transform the value if set
\end{itemize}

\item \textbf{Gating}:
\begin{itemize}
\item Compute innovation $\nu_t = z_t - h(\bar{\mu}_t)$
\item Mahalanobis $d_M^2 = \nu_t^\top S_t^{-1} \nu_t$
\item Compare against \texttt{odom1\_pose\_rejection\_threshold}
\item Reject if too far
\end{itemize}

\item \textbf{Update step}:
\begin{itemize}
\item Kalman gain $K_t = \bar{\Sigma}_t H_t^\top S_t^{-1}$
\item State update $\mu_t = \bar{\mu}_t + K_t \nu_t$
\item Covariance update $\Sigma_t = (I - K_t H_t) \bar{\Sigma}_t$
\end{itemize}

\item \textbf{Publish}:
\begin{itemize}
\item TF if \texttt{publish\_tf}
\item \texttt{/odometry/filtered\_map} (or \texttt{/odometry/filtered})
\item Diagnostics if \texttt{print\_diagnostics}
\end{itemize}
\end{enumerate}

\subsection{The Trust Triangle}

Three covariances determine the filter's behavior:
\begin{itemize}
\item $\Sigma_0$ (initial trust in self)
\item $R_t$ (ongoing trust in motion model)
\item $Q_t$ (trust in each sensor)
\end{itemize}

The Kalman gain magnitude:
\[
K_t \sim \frac{\bar{\Sigma}_t}{\bar{\Sigma}_t + Q_t}
\]
controls measurement weight. The relative scale of $R_t$ vs $Q_t$ determines steady-state behavior:

\begin{center}
\begin{tabular}{ll}
\toprule
\textbf{Setting} & \textbf{Effect} \\
\midrule
$R_t$ large, $Q_t$ small & Trust sensors. Output noisy. Fast response. \\
$R_t$ small, $Q_t$ large & Trust model. Output smooth. Slow response. \\
$R_t \sim Q_t$ & Balanced. Typical for well-tuned filter. \\
$R_t \gg Q_t$ \emph{and} model wrong & Filter chases noisy sensors, output unstable. \\
$R_t \ll Q_t$ \emph{and} model wrong & Filter ignores sensors, drifts away from truth. \\
\bottomrule
\end{tabular}
\end{center}

\subsection{Why EKF \#1 and EKF \#2 Have Different $R_t$}

\begin{intuitionbox}{The Key Insight}
Both EKFs run on the same robot. They differ only in which sensors they fuse and what frame they operate in. Why are their process noise covariances different by $10\times$?

\textbf{Answer:} EKF \#1 (odom) sees only continuous, high-rate sensors (wheel + IMU). Between sensor ticks, very little time passes ($\Delta t = 0.02$ s), so process noise per step is small.

EKF \#2 (map) operates between GPS fixes (10 Hz $\to$ 0.1 s). During those 100 ms, the robot can drift significantly relative to the global frame. To prepare for the GPS update, $\bar{\Sigma}_t$ must grow large enough that $K_t$ is high when GPS arrives. Hence $R_t$ is set higher to inflate $\bar{\Sigma}$ faster.

The math: between fixes,
\[
\bar{\Sigma}_t = \bar{\Sigma}_{t-1} + R\, \Delta t
\]
We want $\bar{\Sigma}_{xx}$ to be comparable to $Q_{xx}^{\text{GPS}} = 0.26$ when GPS arrives. With $\Delta t = 0.1$ and $R_{xx} = 0.5$: $\Sigma$ grows by 0.05 per cycle, reaches ~0.5 in 10 cycles --- well above GPS variance. Kalman gain is high; filter snaps to GPS.

If $R_{xx}$ were 0.05 like EKF \#1, $\Sigma$ would only grow to 0.005 between fixes --- way below GPS variance. Kalman gain $\approx 0.005 / 0.265 \approx 2\%$. Filter would barely move toward GPS. Result: GPS rejected as outlier, map estimate drifts.
\end{intuitionbox}

% ============================================================
\section{Tuning Workflow}
% ============================================================

\subsection{Initial Sanity Check}

\begin{enumerate}
\item All sensors publishing? (\texttt{ros2 topic hz})
\item All TFs resolving? (\texttt{ros2 run tf2\_tools view\_frames})
\item Each sensor's frame\_id correct in URDF?
\item EKF receiving messages? (\texttt{/diagnostics})
\end{enumerate}

\subsection{Stationary Test}

Robot should be still:
\begin{itemize}
\item Output velocity $\approx 0$
\item Output position drifts $\approx 0$ over minutes
\item Output yaw stable to $< 0.01$ rad/min
\end{itemize}

Failures:
\begin{itemize}
\item Velocity drifts $\to$ accelerometer bias not cancelled. Lower $R_{\dot{x}\dot{x}}$, increase $Q$ for accel.
\item Yaw drifts fast $\to$ gyro bias too large. Use \texttt{imu\_filter\_madgwick} for orientation, increase its bias estimation gain.
\end{itemize}

\subsection{Straight-Line Test}

Drive 10 m in straight line. Check:
\begin{itemize}
\item Output position matches (within sensor tolerance)
\item No spurious yaw drift (filter shouldn't ``believe'' robot turned)
\end{itemize}

\subsection{Turn-in-Place Test}

Spin 360$^\circ$ at 0.5 rad/s. Check:
\begin{itemize}
\item Output integrates to $\approx 0$ net displacement
\item Yaw integrates to $\approx 2\pi$
\end{itemize}

Mismatch in yaw indicates baseline calibration error or gyro scale factor.

\subsection{NIS Validation}

Compute Normalised Innovation Squared per measurement:
\[
\mathrm{NIS}_t = \nu_t^\top S_t^{-1} \nu_t
\]
Should be $\chi^2_k$ distributed where $k$ is measurement dimension.

\begin{center}
\begin{tabular}{lll}
\toprule
\textbf{Average NIS} & \textbf{Diagnosis} & \textbf{Fix} \\
\midrule
$\approx k$ & Well-tuned & --- \\
$\gg k$ & Overconfident filter & Increase $R_t$ or $Q_t$ \\
$\ll k$ & Underconfident filter & Decrease $R_t$ or $Q_t$ \\
\bottomrule
\end{tabular}
\end{center}

\subsection{Common Failure Diagnoses}

\begin{warnbox}{Filter Output Stuck}
Symptom: filter ignores GPS, position frozen.

Causes:
\begin{enumerate}
\item Initial covariance too small (\texttt{1e-9} for $x, y$ in EKF \#2) $\to$ Kalman gain near zero. Increase to 1.0+.
\item Rejection threshold too tight $\to$ all GPS rejected. Increase to 5.0.
\item GPS covariance in message too large $\to$ Kalman gain near zero. Override with smaller value.
\end{enumerate}
\end{warnbox}

\begin{warnbox}{Filter Output Jumpy}
Symptom: position bounces around.

Causes:
\begin{enumerate}
\item Rejection threshold too loose $\to$ outliers accepted. Decrease.
\item $R_t$ too high $\to$ filter chases every measurement. Decrease.
\item Sensor noise underestimated in $Q_t$. Increase.
\end{enumerate}
\end{warnbox}

\begin{warnbox}{Yaw Drifts}
Symptom: heading slowly rotates over time.

Causes:
\begin{enumerate}
\item Wheel baseline miscalibrated. Recalibrate.
\item IMU yaw bit not enabled in \texttt{imuN\_config}. Enable.
\item No magnetometer + long run $\to$ inevitable. Use loop closure or GPS heading.
\end{enumerate}
\end{warnbox}

\begin{warnbox}{TF Errors}
Symptom: ``Lookup would require extrapolation into future'' / ``into past''.

Causes:
\begin{enumerate}
\item Clock skew. Sync system clocks.
\item \texttt{transform\_time\_offset} set wrong. Try 0.05.
\item Sensor timestamps in wrong unit. Check ros2 topic echo.
\end{enumerate}
\end{warnbox}

% ============================================================
\section{Quick Reference Tables}
% ============================================================

\subsection{All Parameters at a Glance}

\begin{center}
\footnotesize
\begin{longtable}{p{4.0cm} p{2.0cm} p{2.0cm} p{6.0cm}}
\toprule
\textbf{Parameter} & \textbf{Type} & \textbf{Project value} & \textbf{Effect of increasing} \\
\midrule
\endfirsthead
\multicolumn{4}{l}{\itshape continued from previous page} \\
\toprule
\textbf{Parameter} & \textbf{Type} & \textbf{Project value} & \textbf{Effect of increasing} \\
\midrule
\endhead
\multicolumn{4}{l}{\itshape continued on next page} \\
\endfoot
\bottomrule
\endlastfoot
\texttt{frequency} & Hz & 50.0 & smoother, more CPU \\
\texttt{sensor\_timeout} & s & 0.1 & wait longer for laggy sensors \\
\texttt{two\_d\_mode} & bool & true & lock z, roll, pitch \\
\texttt{world\_frame} & string & odom/map & determines TF role \\
\texttt{publish\_tf} & bool & true & node owns TF link \\
\texttt{use\_control} & bool & true & use cmd\_vel for prediction \\
\texttt{control\_timeout} & s & 0.2 & ignore stale commands \\
\texttt{acceleration\_limits} & m/s² & [1.0, ..., 1.5] & allow sharper accel \\
\texttt{acceleration\_gains} & [0,1] & [0.8, ..., 0.9] & trust commanded accel more \\
\texttt{process\_noise\_covariance} ($R$) & 15$\times$15 & varies & faster filter, noisier \\
\texttt{initial\_estimate\_covariance} ($\Sigma_0$) & 15$\times$15 & varies & first measurement dominates \\
\texttt{odomN\_config} & 15 bool & varies & include more state dims \\
\texttt{odomN\_differential} & bool & false & use deltas, not absolute \\
\texttt{odomN\_relative} & bool & false & first msg = origin \\
\texttt{odomN\_pose\_rejection\_threshold} & $d_M^2$ & 3.0 (GPS) & accept more outliers \\
\texttt{imuN\_remove\_gravitational\_acceleration} & bool & true & subtract gravity \\
\texttt{magnetic\_declination\_radians} & rad & 0.0 & rotate map frame \\
\texttt{yaw\_offset} & rad & 0.0 & align IMU 0 with map +x \\
\texttt{zero\_altitude} & bool & true & ignore GPS altitude \\
\texttt{wait\_for\_datum} & bool & false & require explicit datum \\
\texttt{delay} (navsat) & s & 3.0 & wait longer before locking transform \\
\end{longtable}
\end{center}

\subsection{Project Configuration Summary --- Side-by-Side}

\begin{center}
\footnotesize
\begin{tabular}{p{4.0cm} c c}
\toprule
\textbf{Parameter} & \textbf{EKF \#1 (odom)} & \textbf{EKF \#2 (map, post-debug)} \\
\midrule
\texttt{world\_frame} & odom & map \\
\texttt{frequency} & 50 & 50 \\
\texttt{two\_d\_mode} & true & true \\
\texttt{use\_control} & true & false \\
Wheel odom (config) & $v_x, \omega_z$ & $v_x, v_y, \omega_z$ \\
IMU (config) & orient + $\omega$ + accel & $\omega$ only (no accel) \\
\texttt{imu0\_differential} & false & false \\
Visual odom & yes (pose, diff) & --- (removed) \\
LiDAR odom & yes (pose+yaw, diff) & yes (yaw only, diff mode) \\
GPS & --- & \textbf{yes (RTK, $R\!=\!10^{-4}$)} \\
GPS rejection threshold & --- & 100 \\
$R_{xx}$ & 0.05 & 0.05 \\
$R_{\dot{x}\dot{x}}$ & 0.025 & 0.001 \\
$R_{\psi\psi}$ & 0.06 (or 0.08) & 0.001 \\
$R_{\dot{\psi}\dot{\psi}}$ & 0.04 & 0.001 \\
$\Sigma_0^{xx}$ & \texttt{1e-9} & \texttt{1e-9} \\
$\Sigma_0^{\psi\psi}$ & \texttt{1e-9} & \texttt{1e-9} \\
\midrule
Map x,y drift (stationary) & --- & $\le 3$~cm (RTK noise floor) \\
Map yaw drift (stationary) & --- & $\le 0.1^\circ$ (locked) \\
Global anchoring & --- & GPS (RTK) \\
\bottomrule
\end{tabular}
\end{center}

\textbf{Reading the table}: EKF~\#2 uses a leaner sensor set than EKF~\#1 (no VO, no IMU accel) plus tight velocity/yaw process noise. Three tweaks vs ``canonical'': (1) GPS rejection threshold 100$\sigma$ (RTK noise too low to produce realistic outliers), (2) tight $Q[\psi], Q[\dot{y}], Q[\dot{\psi}]$ (kinematic + cross-cov constraints), (3) navsat node consumes EKF~\#2 output (\emph{not} EKF~\#1) to break the circular reference loop documented in \texttt{ekf\_drift\_debug.pdf}.

% ============================================================
\section*{Closing Notes}
\addcontentsline{toc}{section}{Closing Notes}
% ============================================================

Every parameter in \texttt{robot\_localization} corresponds to a specific term in the EKF equations. The \texttt{config} bits build the observation matrix $H_t$. The \texttt{covariance\_diagonal} overrides shape the measurement noise $Q_t$. The \texttt{process\_noise\_covariance} is the motion noise $R_t$. The \texttt{initial\_estimate\_covariance} is $\Sigma_0$. The \texttt{rejection\_threshold} parameters implement Mahalanobis gating $d_M^2 < \text{threshold}$. The \texttt{world\_frame} parameter determines whether the node plays the local or global role in the dual-EKF design.

Tuning is not arbitrary --- each knob has a mathematical home. When the filter misbehaves, identify which equation governs the symptom, then adjust the parameter that controls that equation. The relative scale of $R$ to $Q$ determines steady-state Kalman gain; the relative scale of $\Sigma_0$ to $Q$ at startup determines initial trust in the first measurement.

\textbf{Update (post-debug, 2026-05-09):} The earlier ``deliberate omission of GPS from EKF~\#2'' has been reversed. GPS is now fused in EKF~\#2 with the following preconditions met:
\begin{itemize}[noitemsep]
\item URDF GPS noise tightened to RTK grade ($\sigma\!\approx\!1$~cm horizontal). Relay covariance matched at $10^{-4}$~m$^2$.
\item LiDAR yaw added in differential mode (\texttt{odom2\_config[5]=true}, \texttt{differential=true}). Provides the absolute yaw observation needed to suppress GPS-driven yaw coupling.
\item Wheel $v_y\!=\!0$ enabled (\texttt{odom0\_config[7]=true}). Enforces non-holonomic constraint at the measurement level.
\item IMU linear acceleration disabled in EKF~\#2. Wheel + GPS provide all needed velocity/position information; accel only injects bias.
\item Process noise tightened: $Q[\psi]=Q[\dot{x}]=Q[\dot{y}]=Q[\dot{\psi}]=0.001$. Non-holonomic constraint + cross-covariance suppression.
\item GPS rejection threshold raised to 100$\sigma$. RTK is too clean to produce realistic outliers; gating only blocks legitimate corrections.
\item \texttt{navsat\_transform\_node} input remapped from EKF~\#1 (\texttt{/odometry/local}, odom frame) to EKF~\#2 (\texttt{/odometry/global}, map frame). Without this remap, GPS in odom frame gets re-transformed by EKF~\#2's own \texttt{map}$\to$\texttt{odom} TF, producing self-consistent garbage that never corrects state. See \texttt{report/ekf\_drift\_debug.pdf} for the 11-step case study.
\end{itemize}

Result: stationary hold at GPS within 1--3~cm, yaw locked within $0.1^\circ$. Original 3.56$°/$s yaw leak observation came from a config with loose $Q[\psi]$ and consumer-grade GPS noise; the same architectural pattern is stable once those underlying issues are fixed.

Combine this document with \texttt{EKF\_Deep\_Dive\_Guide.pdf} (theory) for a complete picture: theory explains \emph{why} the equations look the way they do; this document explains \emph{which knob controls what}, grounded in the actual project YAML.

\end{document}
